\documentclass[a4paper, serif, 10pt]{moderncv}

%% moderncv
% style and color
\moderncvstyle{banking}
\moderncvcolor{black}

%% page layout/margins
\usepackage[top=2.5cm, bottom=2.5cm, left=1.5cm, right=1.5cm]{geometry}

\nopagenumbers{}

%% Babel config for Brazilian Portuguese language
% ref:
% - https://latex3.github.io/babel/guides/locale-portuguese.html
\usepackage[brazilian]{babel}

%% fontspec
\usepackage{fontspec}

\IfFontExistsTF{Linux Libertine}{
  \setmainfont[Ligatures={TeX, Common}, Numbers=OldStyle]{Linux Libertine}
}{}
\IfFontExistsTF{Linux Libertine O}{
  \setmainfont[Ligatures={TeX, Common}, Numbers=OldStyle]{Linux Libertine O}
}{}

%% CTeX
% CJK support, used to show CJK characters in the resume
%
% - fontset=none: disable builtin fontset but instead set the CJK font manually
% - heading=false: disable ctex heading
% - punct=kaiming: use kaiming punctuations styles for CJK
% - scheme=plain: use plain scheme, do not override \normalsize font size
% - space=auto: space settings for CJK characters
%
% ref:
% - http://ctan.mirrorcatalogs.com/language/chinese/ctex/ctex.pdf
\usepackage[UTF8, heading=false, punct=kaiming, scheme=plain, space=auto]{ctex}

\IfFontExistsTF{Noto Serif CJK SC}{
  \setCJKmainfont{Noto Serif CJK SC}
}{}
\IfFontExistsTF{Noto Sans CJK SC}{
  \setCJKsansfont{Noto Sans CJK SC}
}{}

%% line spacing
\usepackage{setspace}
\setstretch{1.125}

%% URL styling - use normal text instead of monospace
\urlstyle{same}

% Auto-underline all links
% Defer until \begin{document} so hyperref is already loaded
\AtBeginDocument{%
  \let\oldhref\href
  \renewcommand{\href}[2]{\oldhref{#1}{\underline{#2}}}%
}

%% Patch \httplink and \httpslink to handle full URLs
% The original moderncv macros always prepend http:// or https:// to URLs.
% This causes issues when the URL already has a protocol (e.g., https://example.com)
% resulting in malformed URLs like https://https://example.com.
% This patch checks if the URL already contains "://" and uses it directly if so.
% Note: We use \str_set:Nx (x-expansion) to expand macros like \@homepage first.
\ExplSyntaxOn
\str_new:N \g_yamlresume_url_str

\RenewDocumentCommand{\httpslink}{O{}m}{
  \str_set:Nx \g_yamlresume_url_str {#2}
  \str_if_in:NnTF \g_yamlresume_url_str {://}
    { \url{#2} }
    { \url{https://#2} }
}

\RenewDocumentCommand{\httplink}{O{}m}{
  \str_set:Nx \g_yamlresume_url_str {#2}
  \str_if_in:NnTF \g_yamlresume_url_str {://}
    { \url{#2} }
    { \url{http://#2} }
}
\ExplSyntaxOff

\name{Ana Beatriz Silva}{}
\title{Gerente de Projetos Ágeis | PMP | Scrum Master}
\phone[mobile]{+55 11 98765-4321}
\email{ana.silva@email.com}
\homepage{https://www.linkedin.com/in/anabeatrizsilva}

\address{Rua das Flores, 123, São Paulo, SP, Brasil, 01234-567}{}{}

\extrainfo{{\small \faLinkedin}\ \href{https://www.linkedin.com/in/anabeatrizsilva}{@anabeatrizsilva}}

\begin{document}

\maketitle

\section{Informações Básicas}

\cvline{}{\begin{itemize}
\item Gerente de Projetos com mais de 8 anos de experiência liderando equipes multidisciplinares em projetos de tecnologia e transformação digital.
\item Especialista em metodologias ágeis (Scrum, Kanban) e tradicionais (PMBOK), com certificação PMP e Scrum Master.
\item Histórico comprovado de entregas dentro do prazo e orçamento, com foco em qualidade e satisfação do cliente.
\item Habilidade em gestão de stakeholders, mitigação de riscos e otimização de processos.
\end{itemize}}
\section{Formação}

\cventry{mar. de 2010–dez. de 2014}
        {Bacharelado, Administração de Empresas, Nota: 8.5}
        {Universidade de São Paulo}
        {\url{https://www.usp.br}}
        {}
        {\begin{itemize}
\item Formação sólida em administração com ênfase em gestão de projetos.
\item Desenvolvimento de habilidades em liderança, comunicação e análise de dados.
\end{itemize}
\textbf{Cursos}: Gestão de Projetos, Planejamento Estratégico, Gestão de Pessoas, Finanças Corporativas, Marketing, Estatística}

\cventry{fev. de 2016–dez. de 2017}
        {Mestrado, Gestão de Projetos, Nota: 9.0}
        {Fundação Getulio Vargas}
        {\url{https://www.fgv.br}}
        {}
        {\begin{itemize}
\item Especialização em gestão de projetos com foco em metodologias ágeis e tradicionais.
\item Trabalho de conclusão sobre transformação ágil em empresas brasileiras.
\end{itemize}
\textbf{Cursos}: Gerenciamento de Projetos Complexos, Métodos Ágeis, Liderança de Equipes}
\section{Experiência Profissional}

\cventry{jan. de 2020–Atual}
        {Gerente de Projetos Sênior}
        {Tech Solutions Ltda}
        {\url{https://www.techsolutions.com.br}}
        {}
        {\begin{itemize}
\item Liderança de uma equipe de 15 pessoas no desenvolvimento de uma plataforma SaaS, resultando em um aumento de 30\% na receita anual.
\item Implementação de metodologia Scrum, reduzindo o time-to-market em 25\%.
\item Gestão de orçamento de R\$ 5 milhões, com entrega dentro do prazo e custo.
\item Interface com stakeholders de alto nível, incluindo C-level, para alinhamento estratégico.
\end{itemize}
\textbf{Palavras-chave}: Scrum, SaaS, Gestão de Orçamento, Liderança}

\cventry{mar. de 2017–dez. de 2019}
        {Gerente de Projetos}
        {Inova Consultoria}
        {\url{https://www.inovaconsultoria.com.br}}
        {}
        {\begin{itemize}
\item Gerenciamento de projetos de consultoria em transformação digital para clientes dos setores financeiro e varejo.
\item Coordenação de equipes remotas e presenciais, garantindo a entrega de soluções personalizadas.
\item Desenvolvimento de cronogramas, planos de comunicação e análise de riscos.
\item Aumento da satisfação do cliente em 20\% através de melhorias nos processos de entrega.
\end{itemize}
\textbf{Palavras-chave}: Transformação Digital, Consultoria, Gestão de Riscos, Comunicação}

\cventry{jan. de 2015–fev. de 2017}
        {Analista de Projetos}
        {Startup XYZ}
        {\url{https://www.startupxyz.com.br}}
        {}
        {\begin{itemize}
\item Suporte ao gerente de projetos na elaboração de cronogramas e acompanhamento de tarefas.
\item Utilização de ferramentas como Jira e Trello para monitoramento de sprints.
\item Participação em reuniões diárias e retrospectivas, contribuindo para a melhoria contínua.
\item Auxílio na documentação de projetos e relatórios de status.
\end{itemize}
\textbf{Palavras-chave}: Jira, Trello, Ágil, Documentação}
\section{Idiomas}

\cvline{Português}{Nativo ou Bilíngue \hfill \textbf{Palavras-chave}: Língua materna}
\cvline{Inglês}{Proficiência Profissional Fluente \hfill \textbf{Palavras-chave}: TOEFL 105}
\cvline{Espanhol}{Proficiência Profissional Limitada}
\section{Competências}

\cvline{Gestão de Projetos}{Especialista \hfill \textbf{Palavras-chave}: PMBOK, Agile, Scrum, Kanban, Waterfall}
\cvline{Liderança}{Avançado \hfill \textbf{Palavras-chave}: Team Management, Mentoring, Negotiation, Conflict Resolution}
\cvline{Ferramentas}{Avançado \hfill \textbf{Palavras-chave}: Jira, Trello, MS Project, Asana, Slack}
\cvline{Metodologias Ágeis}{Especialista \hfill \textbf{Palavras-chave}: Scrum, Kanban, SAFe, Lean}
\section{Prêmios}

\cventry{nov. de 2022}
        {PMI São Paulo}
        {Prêmio de Excelência em Gestão de Projetos}
        {}
        {}
        {Reconhecimento pelo sucesso na entrega de um projeto crítico de transformação digital, superando as metas estabelecidas.}
\section{Certificados}

\cventry{mar. de 2018}
        {Project Management Institute (PMI)}
        {Project Management Professional (PMP)}
        {\url{https://www.pmi.org/certifications/project-management-pmp}}
        {}
        {}

\cventry{jun. de 2016}
        {Scrum Alliance}
        {Certified Scrum Master (CSM)}
        {\url{https://www.scrumalliance.org/get-certified/scrum-master-track/certified-scrummaster}}
        {}
        {}
\section{Publicações}

\cventry{ago. de 2021}
        {Transformação Ágil em Empresas Brasileiras}
        {Revista Mundo PM}
        {\url{https://www.mundopm.com.br/artigo/transformacao-agil}}
        {}
        {\begin{itemize}
\item Artigo sobre os desafios e benefícios da adoção de metodologias ágeis em grandes empresas no Brasil.
\item Análise de casos de sucesso e lições aprendidas.
\end{itemize}}
\section{Projetos}

\cventry{jan. de 2021–dez. de 2021}
        {Projeto de implementação de ERP para uma indústria de médio porte.}
        {Implementação de Sistema ERP}
        {\url{https://www.techsolutions.com.br/projetos/erp}}
        {}
        {\begin{itemize}
\item Liderança da equipe de 10 pessoas na migração de sistemas legados para uma solução ERP integrada.
\item Entrega dentro do prazo e orçamento, com treinamento de 200 usuários.
\end{itemize}
\textbf{Palavras-chave}: ERP, Gestão de Mudanças, Treinamento}
\section{Interesses}

\cvline{Corrida}{Maratona, Meia Maratona}
\cvline{Leitura}{Literatura Brasileira, Gestão, Tecnologia}
\section{Voluntariado}

\cventry{jan. de 2018–dez. de 2020}
        {Gerente de Projetos Voluntária}
        {Projeto Crescer}
        {\url{https://www.projetocrescer.org.br}}
        {}
        {\begin{itemize}
\item Coordenação de projetos sociais voltados para educação de jovens em comunidades carentes.
\item Gestão de equipe de 5 voluntários e captação de recursos.
\item Implementação de metodologias ágeis para otimizar as entregas.
\end{itemize}}
    
\end{document}